\chapter{Spin Geometry and the First Entangled State}

\lectureoverview{Entanglement belongs to systems with parts, but the two-spin problem cannot be read honestly until the one-spin machinery is complete. We therefore begin with spin along an arbitrary axis, derive the transition probability between two directions, interpret measurement as a filter, and identify commutation as the criterion for compatible observables. A counting argument then reveals that two spins possess more pure states than two separately prepared spins can supply. The excess states are entangled, and the simplest example already has definite pair correlations while each electron separately is completely unpolarized.}

\section{Finishing the One-Spin Problem}

Bell's theorem can wait. Entanglement itself begins with two electrons, but the remaining structure of one electron will enter every calculation. Let us therefore finish the two-state system before combining two copies of it.

Choose an ordinary spatial unit vector
\begin{equation}
  \hat n=(n_1,n_2,n_3),
  \qquad
  n_1^2+n_2^2+n_3^2=1.
  \label{eq:l04-unit-pointer}
\end{equation}
The labels \(1,2,3\) may equally be called \(x,y,z\). What matters is that the axes form one consistent right-handed coordinate system.

\begin{classroomqa}
\lecturetimestamp{00:02:33}
\audiencequestion{Is the coordinate triad drawn on the board right-handed?}
\lecturerresponse{The first sketch is left-handed. Reversing the second axis, or relabeling two axes consistently, repairs it. The calculation below uses only a fixed right-handed convention; the accidental orientation of the sketch carries no physics.}
\end{classroomqa}

The spin has three observable components, represented by the Pauli matrices
\begin{equation}
  \begin{aligned}
    \sigma_1&=
    \begin{pmatrix}0&1\\1&0\end{pmatrix},
    &
    \sigma_2&=
    \begin{pmatrix}0&-i\\i&0\end{pmatrix},\\[0.5em]
    \sigma_3&=
    \begin{pmatrix}1&0\\0&-1\end{pmatrix}.
  \end{aligned}
\end{equation}
The component along \(\hat n\) is the ordinary spatial dot product of the pointer with this matrix-valued vector:
\begin{equation}
  \sigma_{\hat n}
  \equiv
  \boldsymbol{\sigma}\cdot\hat n
  =
  n_1\sigma_1+n_2\sigma_2+n_3\sigma_3.
  \label{eq:l04-sigma-dot-n}
\end{equation}
This dot product is in physical three-space. It is not the inner product of two state vectors.

\begin{figure}[htbp]
\centering
\includegraphics[width=0.82\linewidth]{lecture_04_figure_02.png}
\caption{The spatial component along \(\hat n\) becomes a linear combination of three matrix observables.}
\lectureframe{00:05:09}
\end{figure}

Substituting the Pauli matrices gives
\begin{equation}
  \sigma_{\hat n}
  =
  \begin{pmatrix}
    n_3 & n_1-in_2\\
    n_1+in_2 & -n_3
  \end{pmatrix}.
  \label{eq:l04-sigma-n-matrix}
\end{equation}
It is convenient to abbreviate
\begin{equation}
  n_-=n_1-in_2,
  \qquad
  n_+=n_1+in_2=n_-^*,
\end{equation}
so that
\begin{equation}
  \sigma_{\hat n}
  =
  \begin{pmatrix}
    n_3&n_-\\
    n_+&-n_3
  \end{pmatrix}.
\end{equation}
The real diagonal and conjugate off-diagonal entries make the operator Hermitian.

\begin{figure}[htbp]
\centering
\includegraphics[width=0.84\linewidth]{lecture_04_figure_03.png}
\caption{The arbitrary-axis spin observable in explicit two-by-two form.}
\lectureframe{00:07:34}
\end{figure}

Operationally, measuring \(\sigma_{\hat n}\) means orienting a strong magnetic field along \(\hat n\) and observing the spin-dependent response, represented in the lecture by whether a photon is emitted. The matrix and the laboratory question are two descriptions of the same measurement.

Two small identities will carry the rest of the calculation. First,
\begin{equation}
  \sigma_{\hat n}^{\,2}=I.
  \label{eq:l04-sigma-n-square}
\end{equation}
This follows either by multiplying Equation~\eqref{eq:l04-sigma-n-matrix} by itself or by using
\begin{equation}
  \sigma_i^2=I,
  \qquad
  \sigma_i\sigma_j=-\sigma_j\sigma_i
  \quad(i\ne j).
\end{equation}
Every eigenvalue \(\lambda\) of \(\sigma_{\hat n}\) therefore satisfies \(\lambda^2=1\), so the only possible outcomes are
\begin{equation}
  \lambda=\pm1.
\end{equation}
Second,
\begin{equation}
  n_+n_-
  =n_1^2+n_2^2
  =1-n_3^2
  =(1-n_3)(1+n_3).
  \label{eq:l04-nplus-nminus}
\end{equation}
This is not a deep theorem, but it is the identity that keeps the eigenvector algebra under control.

\section{From Preparation to Probability}

Prepare the electron with spin \(+1\) along \(\hat n\), then measure along a different direction \(\hat m\). The central one-spin question is
\begin{equation}
  P(m,+\mid n,+)
  =
  P(\sigma_{\hat m}=+1\mid\sigma_{\hat n}=+1).
  \label{eq:l04-central-probability}
\end{equation}
Symmetry suggests that the answer should depend only on the angle between \(\hat n\) and \(\hat m\), and that interchanging the two directions should not change it. Those expectations are useful checks, not substitutes for the probability postulate.

\begin{classroomqa}
\lecturetimestamp{00:14:03}
\audiencequestion{Does quantum mechanics itself tell us that the three Pauli matrices are spin operators, or is that learned from experiment?}
\lecturerresponse{The identification is part of the accumulated empirical structure of quantum mechanics, not the result of one isolated experiment. Repeated spin experiments reveal probabilistic regularities, and the matrix formalism was built through the work of Einstein, Born, Heisenberg, Dirac, and others to encode those regularities. The postulates are exceptionally successful guesses constrained by data.}
\end{classroomqa}

Let
\begin{equation}
  \sigma_{\hat n}\ket{n,+}=\ket{n,+},
  \qquad
  \sigma_{\hat m}\ket{m,+}=\ket{m,+}.
  \label{eq:l04-nm-eigenstates}
\end{equation}
The symbols \(n,+\) and \(m,+\) are labels inside the kets. They identify the eigenstates whose eigenvalue is \(+1\).

\begin{figure}[htbp]
\centering
\includegraphics[width=0.76\linewidth]{lecture_04_figure_04.png}
\caption{The \(+1\) eigenstates for spin along \(\hat n\) and \(\hat m\).}
\lectureframe{00:17:36}
\end{figure}

The first object to calculate is the probability amplitude
\begin{equation}
  \mathcal A(m,+\mid n,+)=\braket{m,+}{n,+}.
\end{equation}
The probability is its absolute square:
\begin{equation}
  P(m,+\mid n,+)
  =
  \left|\braket{m,+}{n,+}\right|^2.
  \label{eq:l04-born-nm}
\end{equation}
Amplitude first, probability second: that order is the essential quantum step.

\begin{figure}[htbp]
\centering
\includegraphics[width=0.78\linewidth]{lecture_04_figure_05.png}
\caption{The overlap between the prepared \(n\)-state and the tested \(m\)-state is the probability amplitude.}
\lectureframe{00:19:59}
\end{figure}

\subsection{Solving for the Eigenstate}

Write a candidate \(+1\) eigenvector of \(\sigma_{\hat n}\) as
\begin{equation}
  \begin{pmatrix}\alpha\\\beta\end{pmatrix}.
\end{equation}
Normalization can wait, and an overall phase has no physical effect. Away from the south-pole coordinate patch, we may therefore choose \(\alpha=1\) and write \(\beta=z\). The top component of
\begin{equation}
  \begin{pmatrix}
    n_3&n_-\\
    n_+&-n_3
  \end{pmatrix}
  \begin{pmatrix}1\\z\end{pmatrix}
  =
  \begin{pmatrix}1\\z\end{pmatrix}
\end{equation}
gives
\begin{equation}
  n_3+n_-z=1,
  \qquad
  z=\frac{1-n_3}{n_-}.
  \label{eq:l04-z-solution}
\end{equation}
The lower component gives the same condition once Equation~\eqref{eq:l04-nplus-nminus} is used.

The unnormalized vector has squared norm
\begin{align}
  \begin{pmatrix}
    1&\dfrac{1-n_3}{n_+}
  \end{pmatrix}
  \begin{pmatrix}
    1\\[3pt]\dfrac{1-n_3}{n_-}
  \end{pmatrix}
  &=
  1+\frac{(1-n_3)^2}{n_+n_-}
  \notag\\
  &=
  1+\frac{1-n_3}{1+n_3}
  =
  \frac{2}{1+n_3}.
  \label{eq:l04-unnormalized-norm}
\end{align}
Hence a normalized choice is
\begin{equation}
  \ket{n,+}
  =
  \sqrt{\frac{1+n_3}{2}}
  \begin{pmatrix}
    1\\[5pt]\dfrac{1-n_3}{n_-}
  \end{pmatrix}.
  \label{eq:l04-normalized-n-state}
\end{equation}
At \(n_3=-1\), this coordinate expression is singular although the state is not; the eigenstate is simply \(\ket d=(0,1)^T\). The singularity belongs to the phase convention, not to the spin.

Replacing \(n\) by \(m\) and conjugating gives
\begin{equation}
  \bra{m,+}
  =
  \sqrt{\frac{1+m_3}{2}}
  \begin{pmatrix}
    1&
    \dfrac{1-m_3}{m_+}
  \end{pmatrix}.
\end{equation}
The overlap looks unpleasant, and multiplying it by its conjugate looks worse. Repeated use of Equation~\eqref{eq:l04-nplus-nminus} reduces the result to
\begin{equation}
  \begin{aligned}
    P(m,+\mid n,+)
    &=\frac{1+n_1m_1+n_2m_2+n_3m_3}{2}\\
    &=\frac{1+\hat n\cdot\hat m}{2}.
  \end{aligned}
  \label{eq:l04-spin-transition-probability}
\end{equation}
If \(\theta_{nm}\) is the angle between the two spatial directions,
\begin{equation}
  P(m,+\mid n,+)
  =
  \frac{1+\cos\theta_{nm}}{2}
  =
  \cos^2\frac{\theta_{nm}}{2}.
  \label{eq:l04-half-angle-probability}
\end{equation}

\begin{classroomqa}
\lecturetimestamp{00:32:47}
\audiencequestion{Must the overlap vanish because the two states are orthogonal?}
\lecturerresponse{No. The labels \(m,+\) and \(n,+\) refer to different measurement axes, not to different eigenvalues of one fixed observable. If \(m=n\), the two states coincide and their overlap is one. They become orthogonal only when the physical directions are opposite, so that one state is the \(+1\) state and the other is the \(-1\) state of the same spin component.}
\end{classroomqa}

\section{What the Probability Formula Means}

Equation~\eqref{eq:l04-spin-transition-probability} passes three immediate tests:
\begin{align}
  \hat m=\hat n
  &\quad\Longrightarrow\quad P=1,
  \notag\\
  \hat m=-\hat n
  &\quad\Longrightarrow\quad P=0,
  \notag\\
  \hat m\cdot\hat n=0
  &\quad\Longrightarrow\quad P=\frac12.
\end{align}
The answer depends only on the relative angle and is symmetric under \(m\leftrightarrow n\). Rotating the entire apparatus changes no prediction.

This is the remarkable compromise made by spin one-half. Along every axis the observed result remains \(+1\) or \(-1\), never a fractional spin value. Smooth rotational dependence appears instead in the probabilities of those two discrete answers.

\begin{classroomqa}
\lecturetimestamp{00:39:20}
\audiencequestion{How do we know that repeated trials concern the same electron, and how is a probability measured?}
\lecturerresponse{One may confine one electron in a cavity and monitor the conserved charge through its electric flux, reprepare that electron along \(\hat n\), and repeat the measurement along \(\hat m\). Equivalently, one may prepare a large ensemble of electrons identically and measure each once. The long-run ratio of plus and minus outcomes is the same in either procedure, within statistical error.}
\end{classroomqa}

\begin{classroomqa}
\lecturetimestamp{00:41:29}
\audiencequestion{Must all members of the ensemble be measured in one collective experiment?}
\lecturerresponse{No. Each member is prepared by the same protocol and measured separately. A million such trials form one statistical experiment only in the sense that their outcome frequencies estimate the same probability.}
\end{classroomqa}

\begin{classroomqa}
\lecturetimestamp{00:42:23}
\audiencequestion{If one electron is reused, how many photons can it emit?}
\lecturerresponse{There is no fixed lifetime quota. Each cycle includes external work when the magnetic field is switched on. If the spin is in the higher-energy orientation, that supplied energy can leave as a photon during relaxation. The apparatus can therefore repeat the preparation-and-emission cycle indefinitely in principle.}
\end{classroomqa}

The formula can be obtained very quickly after rotational invariance is known: choose \(\hat n\) as the third axis, so \(\ket{n,+}=(1,0)^T\), and read the probability from the upper component of \(\ket{m,+}\). But using that shortcut at the outset would assume the symmetry that the explicit calculation has just demonstrated.

\begin{classroomqa}
\lecturetimestamp{00:43:33}
\audiencequestion{Could a better choice of basis remove the ugly arbitrary-axis eigenvectors?}
\lecturerresponse{Yes, once rotations and unitary transformations are available, one can align a coordinate axis with one of the directions and make the calculation short. At this stage the explicit calculation serves a purpose: it verifies that the chosen Pauli matrices really produce a rotationally invariant answer rather than assuming that fact.}
\end{classroomqa}

\begin{classroomqa}
\lecturetimestamp{00:48:30}
\audiencequestion{What is the physical reason that the answer depends only on the angle?}
\lecturerresponse{The structural answer belongs to the quantum theory of symmetry: spatial rotations are represented by unitary transformations acting on the spin states and observables. That language comes later. For now, the angle dependence is an earned result of the Pauli-matrix calculation.}
\end{classroomqa}

\section{Measurement Is a Filter}

A measurement does more than report a number. If \(\sigma_{\hat m}\) is measured and the result is \(+1\), the electron is left in \(\ket{m,+}\). The apparatus has acted as both detector and state preparer, much as a polarizer tests a photon and leaves the transmitted photon polarized along its own axis.

This state-change rule makes the problem of simultaneous measurement concrete. Two observables can have reproducible definite values together only when there are states that are eigenvectors of both.

\begin{classroomqa}
\lecturetimestamp{00:51:29}
\audiencequestion{Can one measurement determine two different spin components, such as \(\sigma_1\) and \(\sigma_2\), simultaneously?}
\lecturerresponse{No. Their eigenvectors are different. A measurement of one component leaves the electron in an eigenstate of that component, but that state is not an eigenstate of the other. The first measurement therefore destroys the definiteness required by the second.}
\end{classroomqa}

The general finite-dimensional statement is most cleanly made for Hermitian operators.

\begin{theorem}
Two Hermitian operators \(A\) and \(B\) are simultaneously diagonalizable if and only if they commute:
\begin{equation}
  [A,B]\equiv AB-BA=0.
  \label{eq:l04-commutator}
\end{equation}
Equivalently, they admit a complete common orthonormal eigenbasis.
\end{theorem}

For the easy direction, suppose
\begin{equation}
  A\ket{\alpha_i}=\lambda_i\ket{\alpha_i},
  \qquad
  B\ket{\alpha_i}=\mu_i\ket{\alpha_i}
\end{equation}
for a complete basis. Then
\begin{align}
  AB\ket{\alpha_i}
  &=\lambda_i\mu_i\ket{\alpha_i},
  \notag\\
  BA\ket{\alpha_i}
  &=\mu_i\lambda_i\ket{\alpha_i},
\end{align}
so \(AB\) and \(BA\) agree on every basis vector. Hence \(AB=BA\). Conversely, if Hermitian \(A\) and \(B\) commute, \(B\) preserves each eigenspace of \(A\); diagonalizing \(B\) within any degenerate eigenspace produces a common eigenbasis.

\begin{figure}[htbp]
\centering
\includegraphics[width=0.76\linewidth]{lecture_04_figure_06.png}
\caption{Two observables may have the same eigenvectors while assigning different eigenvalues to them.}
\lectureframe{00:55:45}
\end{figure}

\begin{classroomqa}
\lecturetimestamp{00:53:56}
\audiencequestion{Could we apply magnetic fields along \(x\) and \(y\) at once and call that a simultaneous measurement?}
\lecturerresponse{The fields add to one resultant direction. The apparatus then measures spin along that resultant axis, perhaps at \(45^\circ\), rather than measuring \(\sigma_x\) and \(\sigma_y\) separately. The attempted construction produces a third observable, not two compatible readings.}
\end{classroomqa}

\section{Every Pure One-Spin State Has an Axis}

Now reverse the eigenvector problem. Given an arbitrary normalized spinor
\begin{equation}
  \ket\psi=
  \begin{pmatrix}\alpha\\\beta\end{pmatrix},
  \qquad
  |\alpha|^2+|\beta|^2=1,
  \label{eq:l04-general-spinor}
\end{equation}
does some spatial direction \(\hat n\) satisfy
\begin{equation}
  \sigma_{\hat n}\ket\psi=\ket\psi?
  \label{eq:l04-reverse-eigenproblem}
\end{equation}
Yes. The lecture first sees why by counting. The two complex components contain four real numbers. Normalization removes one, and the physically irrelevant overall phase removes another, leaving two. A unit pointer also has two independent real parameters, such as longitude and latitude.

\begin{classroomqa}
\lecturetimestamp{01:03:52}
\audiencequestion{How can the state have only two degrees of freedom when \(\alpha\) and \(\beta\) are both complex?}
\lecturerresponse{They begin with four real components. The normalization condition removes one real degree of freedom, and multiplying the whole state by \(e^{i\theta}\) changes no probability, so the common phase removes another. Two physically meaningful parameters remain, matching the two angles of a unit pointer.}
\end{classroomqa}

\begin{classroomqa}
\lecturetimestamp{01:04:20}
\audiencequestion{Do we also need to solve for the eigenvalue?}
\lecturerresponse{Here the target eigenvalue is fixed to \(+1\). We are searching only for the direction whose up-state is the given spinor. The opposite pointer would describe the same ray as a down-state with eigenvalue \(-1\).}
\end{classroomqa}

The count explains why the correspondence should exist. A direct construction confirms it. Define
\begin{align}
  n_1&=\bra\psi\sigma_1\ket\psi
      =2\operatorname{Re}(\alpha^*\beta),
  \notag\\
  n_2&=\bra\psi\sigma_2\ket\psi
      =2\operatorname{Im}(\alpha^*\beta),
  \notag\\
  n_3&=\bra\psi\sigma_3\ket\psi
      =|\alpha|^2-|\beta|^2.
  \label{eq:l04-bloch-components}
\end{align}
Normalization gives \(n_1^2+n_2^2+n_3^2=1\), and direct substitution gives Equation~\eqref{eq:l04-reverse-eigenproblem}. Thus every pure one-spin state is spin-up along one spatial axis.

There is an operational version of the statement. Prepare many copies of the same unknown pure state. Measure batches along trial directions, adjust the magnet, and search for the direction along which every outcome is up. Finite data determine the direction only within uncertainty, but the ideal state has one exact axis.

Questions are part of this method, not interruptions to it. A brief classroom aside recalls ten-year-olds asking freely about string theory, extra dimensions, and whether surrounding neutrinos cause headaches. The details of a question may be wrong; the willingness to expose the point of confusion is still valuable.

\section{Two Spins and Local Operators}

Before combining two spins, record how the Pauli matrices act on the \(\sigma_3\) basis:
\begin{align}
  \sigma_1\ket u&=\ket d,
  &
  \sigma_1\ket d&=\ket u,
  \notag\\
  \sigma_2\ket u&=i\ket d,
  &
  \sigma_2\ket d&=-i\ket u,
  \notag\\
  \sigma_3\ket u&=\ket u,
  &
  \sigma_3\ket d&=-\ket d.
  \label{eq:l04-pauli-action-table}
\end{align}

\begin{figure}[htbp]
\centering
\includegraphics[width=0.80\linewidth]{lecture_04_figure_07.png}
\caption{The Pauli action table repackages the matrices as rules on the up/down basis.}
\lectureframe{01:09:35}
\end{figure}

Now imagine two electrons held at distinguishable locations so that only their spins vary. The four basis states are
\begin{equation}
  \ket{uu},
  \qquad
  \ket{du},
  \qquad
  \ket{ud},
  \qquad
  \ket{dd}.
  \label{eq:l04-two-spin-basis}
\end{equation}
The first label belongs to electron \(1\), the second to electron \(2\). These are mutually orthonormal because they represent distinct joint outcomes of the two third-component measurements.

\begin{figure}[htbp]
\centering
\includegraphics[width=0.82\linewidth]{lecture_04_figure_08.png}
\caption{The four product-basis labels span the two-spin state space.}
\lectureframe{01:11:40}
\end{figure}

Call the operators on electron \(1\) \(\sigma_i\), and those on electron \(2\) \(\tau_i\). More explicitly,
\begin{equation}
  \sigma_i\equiv\sigma_i\otimes I,
  \qquad
  \tau_i\equiv I\otimes\sigma_i.
\end{equation}
The first set acts only on the first label, while the second set acts only on the second:
\begin{align}
  \sigma_1\ket{ud}&=\ket{dd},
  &
  \sigma_1\ket{dd}&=\ket{ud},
  \notag\\
  \sigma_3\ket{dd}&=-\ket{dd},
  &
  \tau_2\ket{ud}&=-i\ket{uu}.
  \label{eq:l04-local-operator-examples}
\end{align}
Once an operator's action is known on the four basis states, linearity determines its action on every two-spin state.

\begin{classroomqa}
\lecturetimestamp{01:17:47}
\audiencequestion{Why is the factor \(-i\) written outside the ket? Does it act on both electrons?}
\lecturerresponse{The symbols inside the ket are labels for a basis configuration. A complex scalar multiplies the whole state vector, so standard notation places it outside. Its origin here is the action of \(\tau_2\) on the second label; writing it outside does not mean that a separate physical operation was applied to both electrons.}
\end{classroomqa}

\begin{classroomqa}
\lecturetimestamp{01:19:44}
\audiencequestion{Does multiplication by \(i\) represent a rotation of the electron in ordinary space?}
\lecturerresponse{Not by itself. It is complex scalar multiplication in state space. Spatial rotations will later be represented by particular unitary operators; an isolated factor \(i\) in this action table is not a literal real-space turn.}
\end{classroomqa}

\section{Six Parameters, Not Four}

A general pure state of the pair is
\begin{equation}
  \ket\Psi
  =
  c_{uu}\ket{uu}
  +c_{du}\ket{du}
  +c_{ud}\ket{ud}
  +c_{dd}\ket{dd}.
  \label{eq:l04-general-two-spin-state}
\end{equation}
Four complex coefficients give eight real numbers. Normalization removes one and overall phase removes one:
\begin{equation}
  8-1-1=6.
  \label{eq:l04-two-spin-parameter-count}
\end{equation}
A separately prepared pure state of electron \(1\) needs two real parameters, and electron \(2\) needs two more, for only four. The two missing parameters are not an accounting error. They signal pure states of the pair that cannot be assigned separate pure states for the parts.

First identify the states that do come from separate preparations. If
\begin{equation}
  \ket{\psi_1}=\alpha_1\ket u+\beta_1\ket d,
  \qquad
  \ket{\psi_2}=\alpha_2\ket u+\beta_2\ket d,
\end{equation}
then
\begin{align}
  \ket{\psi_1}\otimes\ket{\psi_2}
  &=
  \alpha_1\alpha_2\ket{uu}
  +\alpha_1\beta_2\ket{ud}
  \notag\\
  &\quad
  +\beta_1\alpha_2\ket{du}
  +\beta_1\beta_2\ket{dd}.
  \label{eq:l04-product-expansion}
\end{align}
This is a \emph{product state}. Each electron has its own polarization axis, just as it would if the two magnetic preparations were performed independently.

\begin{figure}[htbp]
\centering
\includegraphics[width=0.84\linewidth]{lecture_04_figure_09.png}
\caption{Multiplying two one-spin states produces a restricted four-term product state.}
\lectureframe{01:25:15}
\end{figure}

If the four coefficients are arranged as
\begin{equation}
  C=
  \begin{pmatrix}
    c_{uu}&c_{ud}\\
    c_{du}&c_{dd}
  \end{pmatrix},
\end{equation}
the state is a product precisely when \(C\) has rank one, equivalently
\begin{equation}
  c_{uu}c_{dd}-c_{ud}c_{du}=0.
  \label{eq:l04-product-condition}
\end{equation}
Most choices of four amplitudes do not satisfy this condition. They are entangled.

\section{The First Entangled State}

The lecture's first example is
\begin{equation}
  \ket{\psi_+}
  =
  \frac{1}{\sqrt2}
  \left(\ket{ud}+\ket{du}\right).
  \label{eq:l04-psi-plus}
\end{equation}
It cannot be factored into one state for electron \(1\) and one for electron \(2\). If a third-axis measurement finds electron \(2\) down, electron \(1\) is up; if it finds electron \(2\) up, electron \(1\) is down. The pair has definite information that neither particle possesses separately.

\begin{figure}[htbp]
\centering
\includegraphics[width=0.82\linewidth]{lecture_04_figure_10.png}
\caption{The first entangled superposition contains one up and one down but assigns neither result to a particle in advance.}
\lectureframe{01:29:30}
\end{figure}

To see what either electron looks like alone, calculate the three local expectation values. For the first component,
\begin{equation}
  \sigma_1\ket{\psi_+}
  =
  \frac{1}{\sqrt2}
  \left(\ket{dd}+\ket{uu}\right),
\end{equation}
which is orthogonal to \(\ket{\psi_+}\). Therefore
\begin{equation}
  \bra{\psi_+}\sigma_1\ket{\psi_+}=0.
\end{equation}
Similarly,
\begin{equation}
  \sigma_2\ket{\psi_+}
  =
  \frac{i}{\sqrt2}
  \left(\ket{dd}-\ket{uu}\right),
\end{equation}
so
\begin{equation}
  \bra{\psi_+}\sigma_2\ket{\psi_+}=0.
\end{equation}
For the third component,
\begin{equation}
  \sigma_3\ket{\psi_+}
  =
  \frac{1}{\sqrt2}
  \left(\ket{ud}-\ket{du}\right),
\end{equation}
and the two diagonal contributions cancel:
\begin{equation}
  \bra{\psi_+}\sigma_3\ket{\psi_+}
  =\frac12(1-1)=0.
\end{equation}
Thus
\begin{equation}
  \begin{aligned}
    \bra{\psi_+}\sigma_i\ket{\psi_+}&=0,\\
    \bra{\psi_+}\tau_i\ket{\psi_+}&=0,
    \qquad i=1,2,3.
  \end{aligned}
  \label{eq:l04-zero-local-spin}
\end{equation}

\begin{figure}[htbp]
\centering
\includegraphics[width=0.84\linewidth]{lecture_04_figure_11.png}
\caption{The transformed state and the cancellation \(1-1=0\) establish vanishing local spin expectation values.}
\lectureframe{01:35:50}
\end{figure}

For any direction,
\begin{equation}
  \bra{\psi_+}\sigma_{\hat n}\ket{\psi_+}=0,
  \qquad
  \bra{\psi_+}\tau_{\hat n}\ket{\psi_+}=0.
  \label{eq:l04-zero-arbitrary-local-spin}
\end{equation}
Since each local measurement has only outcomes \(\pm1\), zero expectation means
\begin{equation}
  P(\sigma_{\hat n}=+1)
  =
  P(\sigma_{\hat n}=-1)
  =
  \frac12,
\end{equation}
and likewise for the second electron. No direction polarizes either electron by itself. This is impossible for a pure one-spin state but perfectly possible for one part of an entangled pure state.

\begin{classroomqa}
\lecturetimestamp{01:38:23}
\audiencequestion{Are the half-half local result and the fact that measuring one electron informs us about the other the same claim?}
\lecturerresponse{No. First ignore electron \(2\): every direction measured on electron \(1\) alone gives a half-half distribution. That establishes the absence of any local polarization. Correlation is a separate joint statement about how the two recorded outcomes are related. Entanglement supplies both features, but one should not use one as a substitute for proving the other.}
\end{classroomqa}

\section{Correlation Before Bell}

Correlation by itself is not uniquely quantum. Mix a penny and a dime, separate them without looking, and inspect one. Before inspection either object has a fifty-fifty chance of being the penny, but finding the penny here tells us that the distant object is the dime. The pair carries definite information even when either part by itself does not.

\begin{classroomqa}
\lecturetimestamp{01:41:45}
\audiencequestion{Does the penny-and-dime inference work only if we know the original pair contained exactly one of each?}
\lecturerresponse{Exactly. The classical inference depends on trusted preparation. Likewise, a quantum correlation prediction applies only when the two-electron system has been prepared in the specified joint state.}
\end{classroomqa}

\begin{classroomqa}
\lecturetimestamp{01:42:09}
\audiencequestion{For \(\ket{\psi_+}\), is the probability that both third-component measurements are up equal to zero? Does that prove the spins are opposite along every axis?}
\lecturerresponse{The probability of \(uu\) in the third-axis basis is zero, and the two third-axis outcomes are opposite. But the second conclusion is too strong for this plus-sign state. Its local spin is zero in every direction, while its joint correlation depends on the direction. Universal opposite outcomes belong to the minus-sign singlet introduced at the end of the lecture.}
\end{classroomqa}

\begin{classroomqa}
\lecturetimestamp{01:43:51}
\audiencequestion{Could the classical pair be imagined as two coins welded so that one face is always opposite the other?}
\lecturerresponse{That is a possible classical mechanism for fixed anti-correlation, provided the shaking does not disturb the linkage. It remains a hidden classical arrangement established at preparation. The analogy is deliberately provisional; Bell's theorem will test where every such classical account fails.}
\end{classroomqa}

\begin{classroomqa}
\lecturetimestamp{01:44:32}
\audiencequestion{Where does the law of large numbers enter if one outcome already determines the partner's outcome?}
\lecturerresponse{Conditional certainty is a statement about one trial once the joint state and one result are given. Repeated trials are still needed to verify experimentally that the promised correlation is reliable and to estimate expectation values and imperfect probabilities. The law of large numbers connects those repeated frequencies to the theoretical distribution.}
\end{classroomqa}

The two observers may be separated by an enormous distance. Looking at one result changes what the local observer knows about the other result, but it does not send a controllable message to the distant observer. At this stage the situation still resembles the separated penny and dime.

\begin{classroomqa}
\lecturetimestamp{01:46:46}
\audiencequestion{Is the separated penny-and-dime picture the right way to think about the quantum pair?}
\lecturerresponse{It is the right first classical comparison, not the final quantum explanation. Bell's theorem concerns the measurable difference between a pair carrying pre-existing classical properties and an entangled quantum pair. That distinction is postponed until the necessary correlations have been calculated.}
\end{classroomqa}

\section{The Singlet and Its Preparation}

At the end of the blackboard calculation, replace the plus sign by a minus:
\begin{equation}
  \ket{\psi_-}
  =
  \frac{1}{\sqrt2}
  \left(\ket{ud}-\ket{du}\right).
  \label{eq:l04-singlet}
\end{equation}
The change leaves the local expectations in Equation~\eqref{eq:l04-zero-local-spin} equal to zero and leaves the third-axis outcomes opposite. It does change the transverse joint correlations. This minus-sign state is the spin singlet, the rotationally invariant state in which the two spins are opposite along every common measurement axis.

\begin{figure}[htbp]
\centering
\includegraphics[width=0.84\linewidth]{lecture_04_figure_12.png}
\caption{The final blackboard change replaces the symmetric superposition by the minus-sign singlet while retaining zero local spin.}
\lectureframe{01:47:25}
\end{figure}

In the simplified physical picture used here, bringing two localized electron magnetic moments close enough to interact lets the higher-energy alignment relax by photon emission. With no preferred external axis, the lower spin state is represented by the rotationally symmetric singlet. The exact interaction and the reason the relative phase must be negative require a further calculation.

\begin{classroomqa}
\lecturetimestamp{01:49:02}
\audiencequestion{Is the practical preparation simply to bring two electrons close together and wait?}
\lecturerresponse{In this idealized model, yes: let the spin-spin interaction act and allow radiative relaxation into the lower state. The electrons need only be close enough, for long enough, that their magnetic interaction is effective. A real preparation protocol must control the orbital state and competing interactions as well.}
\end{classroomqa}

\begin{classroomqa}
\lecturetimestamp{01:49:16}
\audiencequestion{Is this just the Pauli exclusion principle?}
\lecturerresponse{No. The two electrons have been assigned distinguishable spatial locations, so their spin discussion does not yet require antisymmetrizing overlapping orbital wavefunctions. Exclusion becomes relevant when identical electrons occupy overlapping one-particle states; the present entanglement argument already works for separated sites.}
\end{classroomqa}

\begin{classroomqa}
\lecturetimestamp{01:50:31}
\audiencequestion{Would an ordinary electron beam naturally contain such interacting pairs?}
\lecturerresponse{Usually not. A beam is dilute, and typical electrons remain too far apart for appreciable spin-spin interaction during the available flight time. Collisions or a deliberately confining interaction can create entangled pairs, but mere membership in the same beam does not.}
\end{classroomqa}

\begin{classroomqa}
\lecturetimestamp{01:51:45}
\audiencequestion{What distance is required for the magnetic interaction?}
\lecturerresponse{There is no universal sharp cutoff; the decisive quantity is the interaction time scale. The lecture gives a rough microscopic estimate of order \(10^{-11}\,\mathrm{cm}\) for rapid coupling. At larger separation the energy splitting is smaller, the emitted photon has longer wavelength, and relaxation takes much longer. A sufficiently long wait can compensate in principle, while an ordinary beam provides too little time.}
\end{classroomqa}

\begin{classroomqa}
\lecturetimestamp{01:53:10}
\audiencequestion{Why must the singlet contain a minus sign rather than a plus sign?}
\lecturerresponse{That is not obvious from the local half-half statistics, because those statistics are identical for the two relative phases. The answer lies in how the state transforms under rotations and how the interaction selects total spin zero. The derivation belongs to the next lecture.}
\end{classroomqa}

The course has now crossed the intended threshold. A single pure spin always has a definite up-axis. A pure pair need not assign any pure state to either member, yet it can encode exact relations between their outcomes. That extra structure is entanglement; Bell's theorem will determine why no classical penny-and-dime story can reproduce all of it.
